% Contenido del póster: Hau, Huang, Lin, Shan, Sheng y Wei (2024), The Journal of Finance. DOI 10.1111/jofi.13384
% Reemplaza a Ghasemi et al. (2024), que no tiene resumen público y es de acceso cerrado (decisión del 2026-11-12).
% Elaborado con apoyo de IA a partir del resumen oficial del artículo (OpenAlex), 2026-11-12.
\def\DOI{10.1111/jofi.13384}
\def\TituloEs{El crédito fintech y el crecimiento empresarial}
\def\TituloOriginal{FinTech credit and entrepreneurial growth}
\def\Autores{Hau, Huang, Lin, Shan, Sheng y Wei}
\def\Anio{2024}
\def\Revista{The Journal of Finance}

\def\Problema{Las pequeñas empresas suelen quedar fuera del crédito tradicional porque los bancos no pueden evaluar bien su riesgo y les exigen garantías. El crédito fintech, basado en datos y algoritmos, promete superar ese problema, pero hace falta medir si realmente ayuda a crecer a las empresas que lo reciben.}
\def\Objetivo{Medir el efecto del acceso al crédito fintech sobre el crecimiento y el desempeño de pequeños vendedores.}
\def\Metodo{Los autores usan las líneas de crédito automáticas que se ofrecen a los vendedores de la plataforma de comercio minorista en línea de Alibaba. Aprovechan una discontinuidad en el algoritmo que decide el crédito: vendedores muy parecidos a ambos lados del umbral reciben o no la línea, lo que permite comparar sus resultados.}
\def\Resultados{\begin{itemize}
  \item El acceso al crédito fintech aumenta el crecimiento de las ventas y de las transacciones del vendedor.
  \item También mejora la satisfacción de los clientes, medida por las calificaciones del producto, del servicio y del envío.
  \item Los efectos son mayores en los vendedores con más asimetría de información sobre su riesgo crediticio y con menos garantías.
\end{itemize}}
\def\Conclusion{El crédito fintech tiene una ventaja informativa sobre la tecnología crediticia tradicional: llega a las empresas más difíciles de evaluar y con menos garantías, y les permite crecer.}
\def\Relevancia{En el Perú, los proyectos empresariales del financiamiento participativo se financian en el \cifra{colocacionEmp}~\% de los casos y por un monto medio de S/~\cifra{montoMedioEmp} mil. Con los datos públicos de la SMV no puede saberse si esas empresas son las que tienen menos garantías ni si crecieron después de financiarse; medirlo exigiría seguir a cada receptor, como hace este estudio.}
\def\PalabrasClave{crédito fintech, crecimiento empresarial, asimetría de información, garantías, regresión discontinua}
% --- Ficha del estudio (herramientas/posteres/ficha_articulos.py) ---
\def\Ambito{China (vendedores de la plataforma de Alibaba)}
\def\Enfoque{Empírico: regresión discontinua}
\def\FiguraPeru{2}
% --- fin de la ficha ---
